%% sections/03-methodology.tex

\section{Methodology}
\label{sec:methodology}

\subsection{The Arc-Completion Design Template}

The arc-completion pattern is operationalized through a four-field design
template that is completed for every significant NPC before any session
begins. The template is embedded in the NPC's design document (in Marathon,
the \texttt{npcs/<slug>.md} file produced by the \textit{npc-architect}
pipeline skill) and is the primary source of evidence for all corpus analyses
in Section~\ref{sec:evaluation}.

\begin{description}
  \item[\textbf{(1) Emotional position:}]
    What the NPC is carrying --- their unresolved emotional state as of
    the beginning of the adventure. Not their backstory in general but
    the specific weight they are currently bearing. This field is phrased
    in present tense, from the NPC's internal perspective, to distinguish
    it from biographical summary. Example: \textit{``Tomek has been holding
    a garrison together for eleven months without orders, without
    reinforcement, and without permission to retreat. He is not proud of
    this. He is tired.''}

  \item[\textbf{(2) Arc-completion act:}]
    The specific statement, act, or gesture the NPC delivers when the
    firing condition is met. Specific, brief, and explicitly marked as
    not expandable by the DM without compromising the moment. The
    arc-completion text is the full text of the moment: not a summary of
    what happens but the actual words spoken or the act performed,
    described at the precision of a stage direction. Example:
    \textit{``Tomek: `Eleven months.' He does not elaborate. He puts the
    token on the ledge and steps back.''}

  \item[\textbf{(3) Firing condition:}]
    The exact party behavior that triggers the arc-completion. Stated as
    a specific observable action, not an emotional state or DM judgment.
    The condition is specific enough that the DM or AI system can determine
    with certainty whether it has been met. The contrast is critical: ``if
    the party earns Tomek's trust'' is not a firing condition; ``if a party
    member acknowledges the length of Tomek's garrison assignment without
    asking him to justify it'' is. Only observable actions or statements
    qualify; inferred emotional states do not.

  \item[\textbf{(4) DM discipline note:}]
    Explicit instruction specifying what the DM must not do. This field is
    negative by design: it specifies constraints rather than permissions.
    Typical contents: do not deliver any version of this arc-completion
    before the firing condition is met; after delivering the arc-completion
    text, the NPC is done, do not continue speaking for them; do not soften
    the arc-completion to make the NPC more sympathetic if they are morally
    grey. A non-firing note may optionally specify the NPC's default
    behavior if the firing condition is not met during the session (typically:
    the NPC remains at their current emotional position; no lesser version
    of the arc-completion fires as a consolation).
\end{description}

\subsection{Character Logic vs. DM Scripting}

The distinction between \textit{character logic} and \textit{DM scripting}
is the central analytical axis of this paper. It is a distinction about
causal origin, not about content: the same NPC line delivered in the same
scene may be a character-logic firing or a DM-scripted moment depending on
whether the party's behavior preceded and produced it.

\paragraph{Character-logic firing.}
An arc-completion fires by character logic when: (a) the session log records
a party action that maps to the pre-specified firing condition; and (b) the
arc-completion act follows as the NPC's immediate response to that party
action. The DM's role is recognition and delivery, not authorship: they
recognize that the condition has been met and deliver the pre-written text.
No judgment about the party's emotional state or narrative readiness is
required.

\paragraph{DM-scripted moment.}
A DM-scripted moment occurs when the DM independently decides that an NPC
should produce an emotional beat --- because the session needs it, because
the DM senses the party is ready, because the encounter feels flat --- and
delivers NPC dialogue or action without a prior party action meeting the
designed firing condition. In AI-simulated TRPG, this category is structurally
absent: the AI running the session reads the NPC file and fires the arc-completion
only when the condition is met. In the Marathon corpus, zero DM-scripted moments
were logged in sessions that used the npc-architect pipeline.

\paragraph{Observable distinction.}
The distinction is observable in the session log: character-logic firings are
preceded by specific, documented party actions that match the firing condition;
DM-scripted firings are not. This is the coding criterion used throughout
Section~\ref{sec:evaluation}.

\subsection{The Small-Moment Principle}

A consistent finding across the corpus is that arc-completion moments are
smaller than practitioners expect. Not a speech, not a monologue, but one
line, one act, one gesture. The design template enforces this as a
constraint through the DM discipline note (``after delivering the arc-
completion text, the NPC is done''), but the corpus data suggests the
constraint reflects something deeper about how emotional payoffs function
in narrative.

When a party has accumulated two to five sessions of investment in an NPC
relationship, the moment of arc-completion is the release of that accumulated
tension. A compressed arc-completion allows the tension to release
completely: one line, one act, silence. An extended arc-completion
redistributes the tension back into exposition, diluting the release.
This is related to the literary principle that the most powerful moments
in well-crafted narratives are often the most compressed, because the
accumulated context does the work and the moment needs only to be the
hinge~\cite{chekhov1886letter}.

The five exemplar NPCs in Table~\ref{tab:exemplars} illustrate this principle.
Across the five, the longest arc-completion text is Varet Senn's journal entry
at 18 words. Three of the five are ten words or fewer.

\subsection{Five Exemplar NPCs}

Table~\ref{tab:exemplars} presents five representative NPC arc-completion
designs from the corpus, selected to illustrate the range of trigger types,
moral valences, and campaign archetypes.

\begin{table*}[t]
\caption{Five Exemplar NPC Arc-Completion Designs from the Marathon Corpus}
\label{tab:exemplars}
\renewcommand{\arraystretch}{1.3}
\begin{tabular}{@{}p{1.8cm}p{1.4cm}p{4.2cm}p{4.4cm}p{2.5cm}@{}}
\toprule
\textbf{NPC} & \textbf{Campaign} & \textbf{Firing Condition} & \textbf{Arc-Completion Act} & \textbf{Trigger Type} \\
\midrule
Tomek Irr
  & C6 (Military Unit)
  & A party member acknowledges the length of Tomek's garrison assignment
    without asking him to justify it; no follow-up question.
  & ``Eleven months.'' He puts the token on the ledge and steps back.
    [No further speech.]
  & Professional / duty \\
\addlinespace
Maret Stave
  & C3 (Halted Spire)
  & A healer in the party treats a patient in front of Maret using the
    same three-breath preparation Maret uses; Maret observes.
  & She steps forward and adjusts the patient's arm without speaking.
    [Act only. She does not explain.]
  & Professional / craft-witness \\
\addlinespace
Hessa Var
  & C4 (Thorngate)
  & The party returns Hessa's family ledger intact, without reading from it
    aloud or referencing its contents.
  & ``You did not read it. Thank you for not reading it.''
    [10 words. Then she closes her door.]
  & Grief-adjacent \\
\addlinespace
Aldras Bent
  & C5 (Thieves Guild)
  & The party presents the ledger-date from the Coldforge account to Aldras
    without demanding an explanation; they show it and wait.
  & ``The handler was Feln. I do not know what he did with it.''
    [He names the handler unprompted. Does not confirm or deny anything else.]
  & Professional / information-holder \\
\addlinespace
Varet Senn
  & C2 (Conclave)
  & Varet is present in the tribunal chamber when the party reads the
    compact clause aloud; they do not ask Varet to speak.
  & He opens his journal and writes one line. He does not show it to anyone.
    Later: the journal entry reads ``They read it. It still means what it said.''
  & Situationally-triggered \\
\bottomrule
\end{tabular}
\end{table*}

\paragraph{Tomek Irr (Professional / Duty).}
Tomek is a garrison NCO in Campaign 6 who has held his post for eleven months
past his relief date without orders. His emotional position is exhaustion
without complaint. His firing condition requires the party to acknowledge
this fact \textit{without} asking for justification --- a passive recognition
rather than an active question. The arc-completion is non-verbal except for
two words: the number of months spoken aloud for the first time, followed
by a physical act (the token on the ledge) that requires no DM elaboration.
This is a paradigmatic professional-code arc: the NPC does not need the party
to make things better; they need the party to see what was done.

\paragraph{Maret Stave (Professional / Craft-Witness).}
Maret is a healer in Campaign 3 whose arc is tied to whether the party
demonstrates competence in her craft rather than asking about it. The firing
condition is unusual: it requires the party to perform a specific ritual
(the three-breath preparation) that Maret uses herself, observed by Maret.
The arc-completion is wordless: an adjustment gesture that is the healer's
equivalent of a nod. The absence of speech is specified explicitly in the
DM discipline note. This instance illustrates the craft-witness arc type:
NPCs whose completion requires being seen, not spoken to.

\paragraph{Hessa Var (Grief-Adjacent).}
Hessa's arc is tied to a family ledger the party has in custody. The firing
condition specifies what the party must \textit{not} do (read from the ledger
aloud, reference its contents) as much as what they must do (return it intact).
This inverted-condition structure --- where the arc fires on restraint rather
than action --- is documented in seven instances across the corpus. The
arc-completion is ten words followed by a physical act (closing the door)
that terminates the scene without DM narration.

\paragraph{Aldras Bent (Professional / Information-Holder).}
Aldras represents the information-holder variant: an NPC who possesses
information the party needs but will only volunteer it under specific
conditions of professional respect rather than interrogation. The party must
show the ledger-date and wait: no demand, no follow-up, no ``who was the
handler?'' The arc-completion names the handler unprompted, confirming the
information the party needed without Aldras admitting involvement in anything.
The DM discipline note specifies: ``Aldras does not confirm or deny anything
beyond naming the handler. The party asked right; he answered right; the
scene is done.'' This instance from Campaign 5 demonstrates that the arc-
completion pattern holds for morally grey NPCs whose arc is not tied to
grief or duty but to professional code.

\paragraph{Varet Senn (Situationally-Triggered).}
Varet's arc-completion is the most structurally complex in the exemplar set.
The firing condition requires Varet to be present in a specific location
(the tribunal chamber) when a specific event occurs (the party reads the
compact clause aloud), combined with a negative constraint (the party does
not ask Varet to speak). The arc-completion is deferred: Varet writes in
his journal during the scene, and the payoff is only accessible to players
who ask to see the journal later. The DM discipline note specifies that Varet's
journal entry must not be offered voluntarily; it fires as a delayed discovery
only if the party seeks it. This instance illustrates the information-access
variant of the situationally-triggered type, in which the arc-completion act
is a future-accessible record rather than an immediate statement.

\subsection{Coding Protocol}

The corpus analysis in Section~\ref{sec:evaluation} codes each arc-completion
instance against three criteria:

\begin{enumerate}
  \item The NPC delivers their arc-completion act (verbatim or close
    paraphrase within two words of the designed text).
  \item The session log records a party action immediately preceding the
    arc-completion that matches the pre-specified firing condition.
  \item The DM does not provide additional NPC dialogue to explain or
    extend the moment after the arc-completion act.
\end{enumerate}

Instances meeting all three criteria are classified as full arc-completion
fires (character logic). Instances meeting criterion 1 but not criterion 2
are classified as DM-scripted. Instances meeting criterion 2 but not
criterion 3 are classified as partial (arc-completion with extension). In the
full corpus, 52 full instances were logged, 7 partial instances, and zero
DM-scripted instances.
